\subsection{Система {\LaTeX}}

{\LaTeX} (произносится «латех» или «лэйтех»)~--- система подготовки документов, построенная поверх языка {\TeX}, разработанного Дональдом Кнутом в 1978~году \cite{knuth1984texbook}. Лесли Лэмпорт создал {\LaTeX} в 1984~году как набор макросов, облегчающих использование {\TeX}. Сегодня актуальной версией является {\LaTeX}2e, а в разработке находится {\LaTeX}3 \cite{latex-project}.

Среди ключевых особенностей системы можно выделить следующие.
\begin{itemize}
    \item \textbf{Мощный набор математических формул.} {\LaTeX} является стандартом для публикации научных статей с формулами~--- пакеты \texttt{amsmath}, \texttt{amssymb} и \texttt{amsthm} предоставляют богатый инструментарий.
    \item \textbf{Разделение содержания и оформления.} Автор пишет текст с разметкой, а система берёт на себя расстановку переносов, выравнивание и нумерацию.
    \item \textbf{Огромная экосистема.} Репозиторий CTAN содержит более 6000 пакетов для решения задач от вставки кода \cite{minted-ctan} до построения диаграмм.
    \item \textbf{Стабильность.} Документы, написанные 20 лет назад, как правило, и сегодня компилируются корректно.
\end{itemize}

К недостаткам {\LaTeX} относятся: многословный синтаксис, сложная диагностика ошибок, а также необходимость нескольких запусков компилятора для корректного формирования ссылок и оглавления. Рекомендуемое справочное издание по системе~--- книга «The {\LaTeX} Companion» \cite{mittelbach2004latex}.

\subsection{Система Typst}

Typst~--- новая система вёрстки документов, разработанная Лауренцем Мэдье (Laurenz Mädje) и Мартином Хаугом (Martin Haug) \cite{typst-github}. Публичная бета-версия вышла в марте 2023~года \cite{typst-blog-2023}. Система написана на языке Rust, что обеспечивает высокую производительность и надёжность.

Typst использует собственный язык разметки, сочетающий простой синтаксис для текста с полноценным языком программирования для сложной логики. Компилятор работает в инкрементальном режиме: при изменении одной части документа повторно обрабатывается только она, что даёт мгновенный предварительный просмотр \cite{typst-docs}.

Основные преимущества Typst:
\begin{itemize}
    \item \textbf{Читаемый синтаксис.} Заголовок задаётся через \texttt{=~Заголовок}, жирный текст~--- через \texttt{*текст*}, что ближе к Markdown.
    \item \textbf{Инкрементальная компиляция.} Изменение одной страницы не требует перекомпиляции всего документа.
    \item \textbf{Понятные ошибки.} Сообщения об ошибках содержат точное указание строки и понятное описание проблемы.
    \item \textbf{Встроенный язык программирования.} Скриптование доступно прямо в разметке без дополнительных пакетов.
    \item \textbf{Современный менеджер пакетов.} Пакеты подключаются и кешируются автоматически.
\end{itemize}
